% Generated draft: requires mathematical and source review.
\documentclass[11pt]{article}
\usepackage[margin=25mm]{geometry}
\usepackage{fontspec}
\setmainfont{Latin Modern Roman}
\usepackage{amsmath,amssymb,mathtools,unicode-math}
\setmathfont{Latin Modern Math}
\usepackage{enumitem,needspace,xurl,hyperref}
\setlength{\parindent}{0pt}
\setlength{\parskip}{5pt}
\setlength{\emergencystretch}{3em}
\newcommand{\N}{\mathbb{N}}
\newcommand{\Z}{\mathbb{Z}}
\newcommand{\Q}{\mathbb{Q}}
\newcommand{\R}{\mathbb{R}}
\newcommand{\C}{\mathbb{C}}
\newcommand{\F}{\mathbb{F}}
\newcommand{\A}{\mathbb{A}}
\newcommand{\PP}{\mathbb{P}}
\newcommand{\E}{\mathbb{E}}
\newcommand{\eps}{\varepsilon}
\newcommand{\dd}{\,\mathrm{d}}
\begin{document}
\section*{UnsolvedMath ID: 30002034}
\textbf{Source problem code:} catalogue entry
\subsection{Definitions and mathematical statement}
A directed graph is a quadruple \(G=(G^0,G^1,s_G,r_G)\), consisting of sets of vertices and edges and source and range maps \(s_G,r_G:G^1\to G^0\). A vertex \(v\) is regular if \(0<|s_G^{-1}(v)|<\infty\). The phrase ``arbitrary graphs'' is retained: no finiteness, countability, or row-finiteness hypothesis is imposed in this draft.

The complex Leavitt path algebra \(L_{\C}(G)\) is the universal associative complex algebra generated by vertices \(v\in G^0\), edges \(x_e\), and formal ghost edges \(y_e\), for \(e\in G^1\), subject to
\[
\begin{gathered}
 v w=\begin{cases}v,&v=w,\\0,&v\ne w,\end{cases}\\
 s_G(e)x_e=x_e=x_e r_G(e),\qquad
 r_G(e)y_e=y_e=y_e s_G(e),\\
 y_e x_f=\begin{cases}r_G(e),&e=f,\\0,&e\ne f,\end{cases}\\
 v=\sum_{s_G(e)=v}x_e y_e\qquad(v\text{ regular}).
\end{gathered}
\]
No common identity is adjoined by this definition. The last relation is imposed only at regular vertices, not at sinks or infinite emitters.

The graph \(C^*\)-algebra \(C^*(G)\) is the universal complex \(C^*\)-algebra generated by mutually orthogonal projections \(p_v\), for \(v\in G^0\), and partial isometries \(t_e\), for \(e\in G^1\), satisfying
\[
\begin{gathered}
 t_e^*t_f=\begin{cases}p_{r_G(e)},&e=f,\\0,&e\ne f,\end{cases}
 \qquad t_e t_e^*\le p_{s_G(e)},\\
 p_v=\sum_{s_G(e)=v}t_e t_e^*\qquad(v\text{ regular}).
\end{gathered}
\]

For permanent UnsolvedMath problem \(30002034\), source code \(\texttt{catalogue entry}\), the question is whether, for every pair of directed graphs \(E,F\),
\[
\bigl(\exists\text{ a ring isomorphism }L_{\C}(E)\longrightarrow L_{\C}(F)\bigr)
\quad\Longleftrightarrow\quad
\bigl(\exists\text{ a complex }{*}\text{-isomorphism }C^*(E)\longrightarrow C^*(F)\bigr).
\]
Both implications are requested. A ring isomorphism here means an additive and multiplicative bijection; complex linearity and preservation of an involution are not required. The right-hand isomorphism is interpreted in the usual category of complex \(C^*\)-algebras. The question concerns existence of isomorphisms, without asserting that a given isomorphism extends or restricts to one on the other side.

These definitions use standard directed-graph conventions. The supplied formulation does not explicitly specify graph cardinality conventions or the category of the right-hand isomorphism, so those conventions require source review. The supplied registry labels the problem open and its statement \texttt{corrected\_verified}; these are supplied curation claims, not independent verification by this draft.
\subsection{Short English statement}
For arbitrary directed graphs, is the existence of an abstract ring isomorphism between their complex Leavitt path algebras equivalent to the existence of a \({*}\)-isomorphism between their graph \(C^*\)-algebras? Both directions are part of the question. The graph and isomorphism conventions stated above require human source review.
\subsection{Sources}
\begin{itemize}
\item Local UnsolvedMath registry, supplied complete source record: permanent ID \(30002034\), source problem code \texttt{catalogue entry}, title ``Leavitt Path Algebra Isomorphisms and Graph \(C^*\)-Algebras,'' in the Oberwolfach Reports Open Problems collection. No individual proposer is supplied; the recorded proposed year is 2012.
\item Supplied source citation: \emph{Mini-Workshop: Endomorphisms, Semigroups and C*-Algebras of Rings (2013)}, Oberwolfach Reports, pp.~1183--1229, DOI \url{https://doi.org/10.4171/owr/2012/20}. The citation's 2013 label differs from the recorded proposed year and the DOI's 2012 component; this discrepancy is retained unresolved.
\item Supplied literature source, labeled ``Abrams--Tomforde'': \url{https://arxiv.org/abs/0810.2569}.
\item Supplied literature source, labeled ``Cortinas classification survey'': \url{https://doi.org/10.1112/blms.13139}.
\item Supplied literature source, labeled ``Oberwolfach report'': \url{https://ems.press/content/serial-article-files/46391}.
\item Literature assessment dated exactly \texttt{2026-08-21}: the supplied curation labels the problem open and reports classification results in special cases and in diagonal-preserving or star-isomorphism settings, while describing the general ring-isomorphism versus graph-\(C^*\)-isomorphism question as open. This is supplied dated curation, not independently verified; its time-sensitive status and cited claims require independent checking.
\item The supplied statement-verification field says that the mathematical content was checked against the official Oberwolfach report and that notation, context, and question form were normalized without changing the source-level claim. That verification is reported as registry provenance; no source was visited for this draft.
\end{itemize}
\end{document}
